\documentclass[11pt]{letter}
\usepackage[utf8]{inputenc}
\usepackage[letterpaper, margin=1.5in]{geometry}
\usepackage{mathptmx}
\usepackage{microtype}
\setlength{\parskip}{12pt}
\setlength{\parindent}{0pt}

\begin{document}

\begin{letter}{Editors\\Nature}

\vspace*{1cm}

\opening{Dear Editors,}

We are writing to submit our manuscript titled \textit{``SuperARC: An Agnostic Test for Narrow, General, and Super Intelligence Based On the Principles of Causal Recursive Compression and Algorithmic Probability''} for consideration for publication in Nature. We believe that our work makes a significant contribution to the fields of artificial intelligence, cognitive science, and computational theory, and aligns well with Nature's mission to publish groundbreaking interdisciplinary research.

Our study introduces SuperARC, a novel, open-ended test grounded in Kolmogorov-Chaitin complexity and algorithmic probability. This test is designed to quantitatively evaluate the intelligence of frontier Artificial General Intelligence (AGI) and Superintelligence (ASI) models without the risk of benchmark contamination. Unlike existing tests that rely on statistical compression, SuperARC challenges AI models on fundamental intelligence features such as synthesis and model creation for inverse problems. Our findings indicate that current Large Language Models (LLMs) are primarily optimized for human language mastery rather than true comprehension and reasoning, highlighting their limitations and the need for more advanced, theoretically grounded approaches.

The significance of our work lies in its potential to redefine how intelligence is measured in artificial systems. By leveraging algorithmic information theory, SuperARC provides a robust framework for evaluating intelligence that is both systematic and objective. This approach not only addresses the limitations of current evaluation methods but also offers a pathway to developing more sophisticated AI models that can achieve universal intelligence.

We believe that our manuscript is particularly suitable for Nature due to its interdisciplinary nature and the broad interest it holds for researchers in AI, cognitive science, and computational theory. The introduction of SuperARC and its implications for the future of AI research align with Nature's commitment to publishing high-impact, transformative science.

We confirm that there are no conflicts of interest associated with this submission. Additionally, we declare that there is no related work currently in press or submitted elsewhere that could be perceived as overlapping with the content of this manuscript.

Thank you for considering our manuscript for publication in Nature. We look forward to the possibility of contributing to your esteemed journal and to the feedback from the review process.
\vspace{1cm}

\closing{Sincerely,}

\vspace{1cm}

Alberto Hernández-Espinosa, Luan Ozelim, Felipe S. Abrahão, and Hector Zenil

\end{letter}
\end{document}